\section{语言不是智能本质，而是一次跃变}

本章从最容易混淆的概念讲起。访谈里王鹤对智能的定义接近“视情况对环境做出反应的能力”。这个定义故意不把语言放在中心，因为许多生物没有高阶语言，也能在环境中展现智能；视觉、触觉、声音和本体感觉都可以成为智能体感知世界的 sensor。语言的重要性在于，它让人类能够压缩经验、传播知识、形成抽象概念和协作学习，因此是一次跃变，而不是智能的唯一来源。

这个判断直接影响具身智能路线。如果把智能等同于语言，就会倾向于认为 LLM 已经解决了大部分问题；如果把智能理解为环境反应能力，就会意识到机器人必须回到感知、动作和反馈。语言模型可以提供常识和任务拆解，但机器人必须在物理世界里行动，失败会改变环境和成本，这就要求另一套数据与评估方式。

\begin{figure}[H]
\centering
\includegraphics[width=0.94\textwidth]{language-intelligence-jump.png}
\caption{语言不是智能本质：语言是人类智能的一次跃变，而不是智能的唯一入口。自制概念图，依据 00:05:58--00:25:08 对谈内容整理。}
\end{figure}

\begin{knowledgebox}{读图：语言增强智能，但不能替代身体}
左侧的 intelligence 强调环境、反应和反馈；右侧 language 强调压缩、传播和协作。读这张图时不要把两边看成对立，而要看成层级关系：身体和感知提供环境接触，语言让经验跨个体传播。具身智能的问题在于，只有语言侧强还不够，物理侧也必须闭环。
\end{knowledgebox}

\subsection{视觉是强 sensor，不是完整智能}

本节进一步区分视觉和智能。视觉是一种非常强的 sensor，比很多感知模态更早出现在高级动物智能演化中；但纯视觉识别往往是 passive perception，即给一张图判断类别、分割区域或估计位置。识别结束后，环境不会因为你说“这是猫”而变化，也不会给出下一步反馈。这就是互联网视觉任务与具身智能任务的差别。

王鹤指出，计算机视觉领域过去依靠 ImageNet、人脸识别、语义分割等任务积累了巨大能力，但这些任务多是互联网数据与人工标注驱动。卷到一定阶段后，研究者开始寻找下一步：让视觉主动起来，让 agent 采取行动，让行动改变观测，再让新的观测指导下一步。这就引出了 Perception-Action Loop。

\begin{warningbox}{不要把识别能力直接等同于具身能力}
能识别杯子、椅子和房间，不等于能移动到椅子旁、抓住杯子、打开柜门或完成整理货架。具身智能需要动作和环境反馈，而不是只在静态图片上判断标签。
\end{warningbox}

\subsection{本章小结}

语言让人类智能发生跃迁，但具身智能不能只靠语言。机器人需要把视觉、动作和环境反馈连成闭环，这也是计算机视觉研究转向具身智能的核心动力。

